\documentclass[11pt,a4paper]{article}
% Share the established exercise-session typography and palette.
\usepackage[margin=18mm,headheight=24pt]{geometry}
\usepackage[T1]{fontenc}
\usepackage[english]{babel}
\usepackage{lmodern}
\usepackage{microtype}
\usepackage{amsmath,amssymb}
\usepackage{siunitx}
\usepackage{tikz}
\usetikzlibrary{arrows.meta,calc,angles,quotes}
\usepackage{enumitem}
\usepackage{xcolor}
\usepackage{fancyhdr}
\usepackage{lastpage}
\usepackage[hidelinks]{hyperref}

\definecolor{VectorGreen}{HTML}{174A38}
\definecolor{VectorAccent}{HTML}{216343}
\definecolor{VectorMint}{HTML}{E2F0E5}
\definecolor{VectorCoral}{HTML}{C45136}
\definecolor{VectorInk}{HTML}{203B31}
\definecolor{VectorMuted}{HTML}{5F7468}

\color{VectorInk}
\setlength{\parindent}{0pt}
\setlength{\parskip}{0.55em}
\setlist[enumerate]{leftmargin=*,itemsep=0.45em,topsep=0.35em}
\sisetup{per-mode=symbol}

\pagestyle{fancy}
\fancyhf{}
\lhead{\textcolor{VectorMuted}{Vectors · Exercise session}}
\rhead{\textcolor{VectorMuted}{Addition and subtraction}}
\cfoot{\textcolor{VectorMuted}{\thepage\ / \pageref*{LastPage}}}
\renewcommand{\headrulewidth}{0.4pt}
\renewcommand{\headrule}{\hbox to\headwidth{\color{VectorMint}\leaders\hrule height \headrulewidth\hfill}}

\newcommand{\sessiontitle}[2]{%
  {\color{VectorGreen}\fontsize{25}{29}\selectfont\bfseries #1\par}
  \vspace{0.2em}
  {\color{VectorMuted}\large #2\par}
  \vspace{0.55em}
  {\color{VectorAccent}\rule{\textwidth}{1.5pt}}
}

\newcommand{\problemheading}[2]{%
  {\color{VectorAccent}\small\bfseries\MakeUppercase{Problem #1}}\par
  {\color{VectorGreen}\LARGE\bfseries #2}\par
  \vspace{0.25em}
}

\newcommand{\prediction}[1]{%
  \begin{center}
  \colorbox{VectorMint}{\parbox{0.92\linewidth}{\textbf{Predict before calculating.} #1}}
  \end{center}
}

\newcommand{\checkpoint}[1]{%
  \begin{center}
  \fcolorbox{VectorAccent}{white}{\parbox{0.91\linewidth}{\textbf{Reasoning checkpoint.} #1}}
  \end{center}
}

\newcommand{\worklines}[1]{%
  \par\vspace{0.25em}
  {\color{VectorMuted}\small\textit{Working and explanation}}\par
  \foreach \n in {1,...,#1}{\vspace{0.54cm}\noindent\textcolor{VectorMint}{\rule{\linewidth}{0.35pt}}\par}
}

\newcommand{\answerbox}[1]{%
  \begin{center}
  \colorbox{VectorMint}{\parbox{0.92\linewidth}{\textbf{Answer.} #1}}
  \end{center}
}

\lhead{\textcolor{VectorMuted}{Dynamics · Exercise session 3}}
\rhead{\textcolor{VectorMuted}{Dot products · No coordinates}}

\rhead{\textcolor{VectorMuted}{Worked solutions}}

\begin{document}
\sessiontitle{Solutions: Let the dot product see}{Instructor copy · Exercise session 3}

\problemheading{1}{A familiar flight, a different route}

Write $c=\vec u\cdot\vec g=-ug\sin\alpha$. Uniform acceleration gives
$\vec r=\vec u t+\tfrac12\vec g t^2$. Take its dot product with itself:
\[
S(t)=\vec r\cdot\vec r=u^2t^2+ct^3+\tfrac14g^2t^4,
\qquad S'(t)=2u^2t+3ct^2+g^2t^3.
\]
Height is the projection opposite to gravity: $\vec g\cdot\vec r=-gh$.
The combination $2S-tS'$ eliminates the $u^2t^2$ term and gives
\[
2S-tS'=-ct^3-\tfrac12g^2t^4
=-t^2\vec g\cdot\vec r=gt^2h.
\]
Thus height can be recovered from squared position and its first derivative:
\[
\boxed{h(t)=\frac{2S-tS'}{gt^2}
=-\frac cg t-\frac12gt^2}\qquad(t>0).
\]
No horizontal or vertical component equations have been used.

\textbf{a) Maximum height.} Complete the square in the recovered height:
\[
h(t)=\frac{c^2}{2g^3}-\frac g2\left(t+\frac{c}{g^2}\right)^2.
\]
Therefore
\[
\boxed{H=\frac{c^2}{2g^3}=\frac{u^2\sin^2\alpha}{2g},
\qquad t_{\mathrm{top}}=-\frac{c}{g^2}=\frac{u\sin\alpha}{g}.}
\]

\textbf{b) Total flight time.} Set the recovered height to zero and take
the nonzero root:
\[
\boxed{T=-\frac{2c}{g^2}=\frac{2u\sin\alpha}{g}=2t_{\mathrm{top}}.}
\]

\textbf{Optional reflection.} Substitute $t_{\mathrm{top}}=-c/g^2$ into
$S'$ to obtain
\[
S'(t_{\mathrm{top}})=2t_{\mathrm{top}}u^2\cos^2\alpha>0.
\]
The projectile is still moving horizontally away from the launch point.
Maximum height is not maximum distance from the origin. A strictly vertical
throw, excluded in the question, would give $S'=0$ at the top.

\checkpoint{A squared position contains information about height, but its stationary points answer a different geometric question.}

\clearpage
\problemheading{2}{Three speeds, an unseen force}

\textbf{a) Recover scalar information without choosing axes.}
Let $A=\vec a\cdot\vec a$ and $B=\vec v_0\cdot\vec a$.
Taking the dot product of $\vec v=\vec v_0+\vec a t$ with itself gives
\[
v^2=v_0^2+2Bt+At^2.
\]
The three readings give, with SI units understood,
\[
v_0^2=25,\qquad 9=25+2B+A,\qquad 25=25+4B+4A.
\]
Thus $A=16$ and $B=-16$, so
\[
\boxed{a=\SI{4}{m/s^2},\qquad |\vec F|=ma=\SI{8}{N}.}
\]
If $\phi$ is the angle between the initial velocity and the force,
\[
\boxed{\cos\phi=\frac{\vec v_0\cdot\vec a}{v_0a}
=-\frac45,\qquad \phi\approx126.9^\circ.}
\]
The absolute direction is undetermined: rotating the entire motion across
the table leaves every speed reading unchanged.

\textbf{b) Complete the square.} With $t$ in seconds,
\[
v^2=25-32t+16t^2=9+16(t-1)^2.
\]
Therefore
\[
\boxed{t_{\min}=\SI{1}{s},\qquad v_{\min}=\SI{3}{m/s}.}
\]
The puck never stops. At the minimum, $\vec v\cdot\vec a=B+At=0$:
velocity and acceleration are perpendicular, and neither is zero.

\textbf{c) Recover a distance without a trajectory.} Write
$\tau=\SI{1}{s}$. Constant acceleration gives
\[
\Delta\vec r=2\tau\vec v_0+\tfrac12\vec a(2\tau)^2
=2\tau(\vec v_0+\vec a\tau)=2\tau\vec v(\tau).
\]
Taking its squared length gives
\[
\boxed{|\Delta\vec r|=2\tau|\vec v(\tau)|=\SI{6}{m}.}
\]
Since $\vec v(\tau)\cdot\vec a=0$, also
$\Delta\vec r\cdot\vec F=0$: the displacement is perpendicular to the
constant force. This is the separation of the endpoints, not the length
of the curved path.

\checkpoint{The readings hide the vector directions, but their squared lengths retain enough information to recover a force magnitude and an angle.}

\clearpage
\problemheading{3}{What force turned it back?}
{\color{VectorCoral}\bfseries Optional · Modeling challenge}

\textbf{a) Let the dot product reveal the force.}
For radial motion, $\vec v=\dot r\,\hat{\vec r}$ and
$\vec F=F(r)\hat{\vec r}$. Dot Newton's law with velocity:
\[
\frac m2\frac{d(v^2)}{dt}=\vec F\cdot\vec v=F(r)\dot r.
\]
On either branch where $\dot r\ne0$, the chain rule gives
\[
\boxed{F(r)=\frac m2\frac{d(v^2)}{dr}.}
\]
The positive slope means repulsion, even on approach, when $\dot r<0$.
The result extends to the turning point by continuity for a smooth force.

\textbf{b) Infer the model, then the force.}
Write $u^2-v^2=D/r^p$. The deficits at $\ell$, $2\ell$, and $4\ell$
are $u^2$, $u^2/8$, and $u^2/64$. Doubling separation divides the
deficit by eight, so $2^p=8$, $p=3$, and $D=u^2\ell^3$. Hence
\[
\boxed{v^2(r)=u^2\left(1-\frac{\ell^3}{r^3}\right),
\qquad F(r)=\frac{3mu^2\ell^3}{2r^4}=\frac{C}{r^4}.}
\]
Here $C=3mu^2\ell^3/2$ is inferred from the original motion.
Doubling separation divides the force by \textbf{sixteen}.
The readings determine this model within the assumed power-law family;
finitely many readings do not determine an arbitrary function uniquely.

\textbf{c) Predict a new turning point.}
Keep $C$ fixed: changing the incoming speed does not change the source.
Integrating $d(v^2)/dr=2C/(mr^4)$ with incoming speed $w$ at infinity gives
\[
v^2=w^2-\frac{2C}{3mr^3}.
\]
At closest approach $v=0$. With $w=2u$,
\[
\boxed{r_{\min}^3=\frac{2C}{3m(2u)^2}=\frac{\ell^3}{4},
\qquad r_{\min}=2^{-2/3}\ell\approx0.630\ell.}
\]
This is an \textbf{extrapolation} below the smallest tabulated separation.
It assumes the inferred force law continues to hold there. The original
expression for $v^2$ applies only to the original incoming speed.

\textbf{d) Zero power does not imply zero force.}
At a turning point $\vec v=\vec0$, so $d(v^2)/dt=2\vec v\cdot\vec a=0$
even when acceleration is nonzero. For the original approach,
\[
\boxed{F(\ell)=\frac{3mu^2}{2\ell}>0.}
\]
The nonzero repulsion reverses the motion. Dividing the power equation by
$\dot r$ exactly at the turn would be invalid; the force follows from the
model and its continuous limit.

\checkpoint{Motion can reveal a force law. A new prediction tests whether that law extends beyond the data used to infer it.}
\end{document}
